\iffalse
集合实际的代码，专注核心功能及其逻辑实现，帮我写我的毕业硕士论文的Experiments部分，用中文写，注意要承接methods部分（参考：），做到逻辑通顺合理，自说自话也要讲好故事，整个experiments包含以下部分：
1.概述，对下面将要描述的方法的概括总结，展现其逻辑过程。
2.我们如何训练yolov11
参考：（1）训练代码，得到yolo11mpre2.pt
（2）微调代码,得到fintunefreeze10.pt
3.我们如何训练REID
    我们使用自己创建的数据集训练了一个ReID 模型，其架构采用MobileNet-B0，考虑到我们检测目标的大小分布从10x10到120x140,而大部分的目标尺寸落在30x30-90x90这个范围，我们将我们的模型改为，输入维度为96x96.参考：
4.DeepSORT中的8维基准方案，参考
5.我们的12维加速度方案，参考
6.我们的5维ekf方案，参考
\fi


\iffalse
\section{Experiments}\label{section::Experiments}
\subsection{概述}
本章围绕“检测—跟踪—计数”的整体流程，系统评估了所提出的人群头部检测与站台场景多目标跟踪计数方法。我们以 \texttt{YOLO11m} 作为检测器，结合基于 MobileNet-B0 的 ReID 外观描述子，并在 DeepSORT 框架内分别接入三类运动模型以开展消融：\(i\) 8维匀速基准模型（标准 DeepSORT 卡尔曼滤波器）；\(ii\) 12维带加速度模型；\(iii\) 结合针孔相机几何约束的 5维物理模型（深度与像素高度成反比）。评测在 \S\ref{section::data} 描述的数据基础上进行，其中wo'me检测器训练与微调所用静态图像数据与多段 MOT-TPCH 视频序列相互补充，既覆盖通用拥挤场景，也覆盖列车进站期间的高密度平台场景。

\subsection{实验设置}
\subsubsection{数据与评测协议}
\begin{itemize}
    \item \textbf{检测器训练/微调：} 第一阶段使用 CrowdHuman 训练 \texttt{YOLO11m}，第二阶段使用我们整合并标注的 Open Sensor Data for Rail 2023、RailEye3D 与自建 TrainPlatformCrowd 三个来源的 2,000 张头部标注图像进行微调，得到 \texttt{YOLO11freeze10.pt}（参见 \S\ref{section::data}）。
    \item \textbf{跟踪与计数评测：} 在自建 MOT-TPCH 数据集的 27 个视频片段上进行评测。该数据集按 \S\ref{section::data} 所述经由时间/空间裁剪与去音频处理，分辨率介于 1920\(\times\)1080 与 2304\(\times\)1296/1536\(\times\)864 之间，并标注连续头部轨迹（MOT 格式）。
    \item \textbf{训练/验证划分：} 检测器的第一阶段与第二阶段按照 \S\ref{section::data} 中给定的训练—验证拆分实施。ReID 模型在 MOT-TPCH 的行人/头部轨迹片段上裁剪生成的实例上训练，并使用留出序列验证其区分度（不与评测序列重叠）。
\end{itemize}

\subsubsection{实现细节}
\paragraph{检测器（YOLO11m）。} 按 \S\ref{section::Methods} 所述，我们选用 \texttt{YOLO11m} 作为精度—速度折中版本。第一阶段在 CrowdHuman 训练获得通用头部检测能力，第二阶段以 \texttt{backbone=10} 冻结微调策略在平台相关图像上进行领域适配，得到 \texttt{YOLO11freeze10.pt} 作为最终检测器权重。跟踪阶段默认使用置信度阈值 \(\geq 0.3\) 与 NMS（与仓库评估脚本一致）。

\paragraph{ReID（MobileNet-B0）。} 我们将 DeepSORT 的外观嵌入网络替换为 MobileNet-B0 结构，考虑平台场景中头部目标的像素高度主要分布在 \(30\times 30\) 到 \(90\times 90\) 区间，采用 \(96\times 96\) 的输入尺寸以兼顾小目标判别力与推理开销。训练数据来自 MOT-TPCH 的轨迹内裁剪实例，损失与优化细节遵循行人重识别常规范式（度量学习/分类损失结合）。

\paragraph{DeepSORT 配置。} 除非另行说明，外观库大小 \texttt{nn\_budget=100}，外观相似度门限 \texttt{max\_cosine\_distance=0.2}，检测框最小高度限制关闭。评估脚本参考仓库中的 \texttt{evaluate\_motchallenge.py}，并对每个序列单独运行与保存结果。

\subsection{三类运动模型}
为验证运动建模对平台场景多目标跟踪稳定性与 ID 一致性的影响，我们在相同检测与 ReID 设置下，对比以下三种卡尔曼滤波方案（接口均与 DeepSORT 兼容）：

\paragraph{8维匀速基准模型（CV-8D）。} 状态向量 \([x, y, a, h, v_x, v_y, v_a, v_h]^\top\)，测量为 \([x, y, a, h]^\top\)。该模型假设常速度运动，使用标准线性观测。其实现对应仓库中的 \texttt{deep\_sort/kalman\_filter.py} 的基类实现，作为广泛采用的基准。

\paragraph{12维加速度模型（CA-12D）。} 在基准模型上扩展加速度分量，状态为 \([x, y, a, h, v_x, v_y, v_a, v_h, a_x, a_y, a_a, a_h]^\top\)，测量仍为 \([x, y, a, h]^\top\)。相较 CV-8D，能更好刻画进站减速、短时加减速与遮挡恢复阶段的状态变化。实现参考 \texttt{deep\_sort/kalman\_filter.py} 中的 \texttt{AccelerationKalmanFilter}（提供速度/加速度平滑与缺帧保守预测）。

\paragraph{5维物理模型（Phys-5D）。} 融合针孔相机几何约束，将像素高度 \(h\) 与深度 \(Z\) 关联：\(Z = fH/h\)。状态取 \([u, v, h, Z, \dot Z]^\top\)，测量为 \([u, v, h]^\top\)。通过 \((X,Y,Z) \leftrightarrow (u,v,h)\) 的投影/反投影在预测时显式更新 \(Z\) 与 \(h\)，并对 \(\dot Z\) 引入减速度约束以贴合列车相对运动的物理先验。实现参考 \texttt{physical3dkf\_5d.py} 与 \texttt{deep\_sort/custom\_kf5d.py}。

\subsection{评估指标}
\paragraph{检测指标。} 在静态图像验证集上报告头部检测的 \textbf{mAP@0.5} 与 \textbf{Recall}，用于衡量微调对平台场景的适配收益。

\paragraph{跟踪指标。} 在 MOT-TPCH 上采用 MOTChallenge 常用指标：\textbf{IDF1}、\textbf{MOTA}、\textbf{ID Switches (IDs)}、\textbf{MT/ML} 等；并给出 \textbf{FPS} 以反映实时性。

\paragraph{计数指标。} 按帧统计检测—跟踪得到的目标数，与人工标注人数计算 \textbf{MAE}、\textbf{RMSE} 与 \textbf{MAPE}。若关注流量（进/出）可增设虚拟计数线，计算事件级精度/召回率与 F1。



\subsection{定量结果}
\paragraph{检测微调效果。} 我们首先验证二阶段微调对头部检测的收益：
\begin{table}[H]
    \centering
    \caption{YOLO11m 头部检测在通用验证集与平台相关验证集上的性能（占位）。}
    \label{tab:detector}
    \begin{tabular}{lccc}
        \toprule
        设置 & 数据集 & mAP@0.5 & Recall \\
        \midrule
        第一阶段（CrowdHuman） & 通用验证集 & [\,] & [\,] \\
        第二阶段微调（本工作） & 平台验证集 & [\,] & [\,] \\
        \bottomrule
    \end{tabular}
\end{table}

\paragraph{跟踪主结果与消融。} 在相同检测与 ReID 条件下，比较三种运动模型：
\begin{table}[H]
    \centering
    \caption{MOT-TPCH 上不同运动模型的跟踪性能比较（占位）。}
    \label{tab:mot-ablation}
    \begin{tabular}{lccccc}
        \toprule
        运动模型 & IDF1(\%) & MOTA(\%) & IDs $\downarrow$ & MT(\%) & FPS \\
        \midrule
        CV-8D（基准） & [\,] & [\,] & [\,] & [\,] & [\,] \\
        CA-12D（加速度） & [\,] & [\,] & [\,] & [\,] & [\,] \\
        Phys-5D（物理） & [\,] & [\,] & [\,] & [\,] & [\,] \\
        \bottomrule
    \end{tabular}
\end{table}

\paragraph{计数精度。} 在有逐帧人数标注的片段上报告计数误差：
\begin{table}[H]
    \centering
    \caption{计数评估（逐帧人数）在 MOT-TPCH 上的结果（占位）。}
    \label{tab:counting}
    \begin{tabular}{lccc}
        \toprule
        方法 & MAE $\downarrow$ & RMSE $\downarrow$ & MAPE(\%) $\downarrow$ \\
        \midrule
        直接检测计数（本工作） & [\,] & [\,] & [\,] \\
        追踪去重计数（本工作） & [\,] & [\,] & [\,] \\
        \bottomrule
    \end{tabular}
\end{table}

\subsection{定性分析}
图示展示典型序列的跟踪效果、遮挡恢复与计数曲线（示例占位）：
\begin{figure}[H]
    \centering
    \includegraphics[width=0.95\linewidth]{pictures/qualitative_tracking_placeholder.png}
    \caption{典型序列跟踪与计数可视化示例（占位）。}
    \label{fig:qual}
\end{figure}

\subsection{可复现性与运行说明}
评测可基于 \texttt{CAmodel/deep\_sort} 仓库脚本复现：
\begin{itemize}
    \item 使用我们训练的检测器权重 \texttt{YOLO11freeze10.pt} 生成每帧检测（头部类别），并按 DeepSORT 所需格式准备特征/检测文件。
    \item 运行 \texttt{evaluate\_motchallenge.py} 对 \texttt{MOT-TPCH} 各序列进行跟踪评估，常用参数：\texttt{--min\_confidence=0.3}，\texttt{--max\_cosine\_distance=0.2}，\texttt{--nn\_budget=100}。
    \item 将 \texttt{kalman\_filter} 替换为上述三种方案中的任意一个以完成消融对比（接口已适配）。
\end{itemize}

为便于后续填充实测数值，本章表格均已预留占位符并给出标签以支持交叉引用。
\fi